\chapter{107}

\section{Bettmeralp/VS (CH), Dienstag, 18.
September}



«Ja, Maman. Du kannst mir glauben, es war verdammt schwer, dich nicht
von Anfang an einzuweihen. Vielleicht verstehst du mich jetzt?»

Louis biss die Lippen zusammen. Zwei Mal hatte sie ihn unterbrochen.
Ansonsten seinen Ausführungen stumm gelauscht. Der Druck in seinem Kopf
wurde weniger, selbst wenn ihn der raue Klang ihrer Stimme
verunsicherte. War sie wütend oder eben erst aufgestanden? Er stellte
sich vor, wie seine Mutter in der Wohnung umherging und wünschte sich,
ihr mit seinen Erklärungen Erlösung aufs Gesicht zu zaubern. Wie damals
und immer, wenn sich die Dinge zwischen ihnen klärten.

Auf ihrer Seite blieb es still. Die Leitung übertrug nur mehr ihre
unregelmässigen Schritte. Er sah sie vor sich, wie sie die halboffenen
Pantoffeln bei jedem Schritt nachzog.

«Maman, ich, ich weiss, dass du gelitten hast und es tut mir unglaublich
leid.»

«Ja, Louis.»

Jetzt brach ein leises Weinen aus ihr heraus.

«Ich, ich musste wieder zurückfliegen. Es gibt noch ein paar Dinge zu
klären, bevor ich zurückkommen kann. Es wird bald sein, Maman.»

Louis wartete.

«Maman, wie, wie geht es dir jetzt?»

Schon beim Aussprechen der Worte wusste er, die Frage war deplatziert.
Reine Verlegenheit. Seine Mutter musste nach wie vor in der Wohnung
umhergehen. Er hörte ihre leisen Schritt. Er wartete.

«Ach, Louis, ich, ich bin ja froh, klärt sich das, das,» sie versuchte
offenbar, sich zusammenzureissen und die Tränen hinunterzuschlucken,
«das Ganze. Das Unglaubliche ist,» sie atmete schwer, «ist nicht
alleine, was dir zugestossen ist, es ist,» wieder pausierte sie. Ihr
gesammelter Schmerz hatte wohl endlich eine offene Tür gefunden und
würde wohl bald mit Wucht aus ihr hinausbrechen. Sie schluchzte.

Louis hatte mit einer starken Reaktion gerechnet, doch ihre heftige
Trauer übertraf seine Ahnung. Seine Beine gaben nach und er liess sich
aufs Bett fallen.

«Maman?»

«Och, Louis,» er hörte sie schlucken, «ich bin gerade überfordert. Weisst
du, ehm, ich kann ja verstehen, was dir passiert ist, nur, ich weiss nicht,
wie klar dir ist, dass auch ich selbst mitten in dieser Geschichte drin
war.»

«Wie meinst du das?»

«Von allen Seiten wurde ich auf dich angesprochen und fühlte mich völlig
hilflos, weil ich nicht verstand, was überhaupt los war und,» jetzt war
ihm, als zähme sie ihre Bitterkeit, «ich weiss gerade nicht, ob es mich
wütend oder einfach masslos traurig macht, Louis. Du wolltest mich
schonen? Du hast mir nicht zugetraut, dass ich mit der Situation umgehen
kann und gerade das, Louis, erschüttert mich.»
Er schluckte. Setzte sich auf. Seine Hände zitterten. Er kämpfte.
"Ich, so habe ich es noch nie angeschaut, Maman."
Für einen Moment blieb es auf beiden Seiten ruhig.
"Ich, ich scheine dich sehr verletzt zu haben, Maman," er brauchte eine Pause, " du hast mit allem recht. Ich
verstehe dich, wirklich. Ich habe dich und alle andern enttäuscht. Ich, ich kann nur hoffen, dass, dass du mir
vergibst.»
«Ach, Louis, was redest du von Vergebung, du bist mein Sohn," sie stöhnte hörbar auf, 
"ich bin erleichtert, dass ich dir meine Seite erklären konnte. Aber, was, was hast du nun vor?»
«Danke, Maman, du kannst dir nicht vorstellen, wie gern ich dich jetzt
in den Arm nehmen möchte,» Louis suchte nach Worten, «nachdem du nun die ganze
Geschichte kennst, kümmere ich mich um all die andern, die ich
enttäuscht habe, um Lucia, Marco, meine Freunde, auch um Grandpère, ihn
habe ich arg vernachlässigt, und,» er klang, als müsse er sich
überwinden, «und sobald ich mich wieder ein bisschen sicherer fühle,
möchte ich Giorgia sehen.»
«Grandpère wird sich freuen. Auch er hat sich grosse Sorgen gemacht. In
seinem Alter! Ich habe versucht, das Schlimmste von ihm fern zu halten,
ihn zu beruhigen.»

Louis sah sich augenblicklich seinem Grossvater gegenüber und sein Herz
machte einen Sprung.

«Und ja, Giorgia, das arme Mädchen. Keiner hier will wahrhaben, was der
Unfall mit ihr angestellt hat. Was, was versprichst du dir von einem
Treffen?»

«Wie ich dir sagte, habe ich sie abstürzen sehen, und, und ich, ich
weiss nicht, aber ich glaube, ich möchte während ihrer Heilungszeit für
sie da sein.»

«Ach ja?»

«Sicher bin ich noch nicht, ich werde darüber nachdenken. Mein
Arbeitgeber Johannes ist so was wie ein Mentor geworden für mich, ein
grossartiger Mensch. Er hat mir angeboten, die nächste Zeit für mich da
zu sein. Ich kann sehr gut mit ihm reden. Ich vertraue ihm. Bevor ich
zurückkomme, werde ich meine diversen Baustellen mit ihm besprechen
können.»

«Baustellen. Ach, Louis.»

Sie seufzte.

«Ja, Maman, der Stoff wird mir nicht ausgehen, mein Gott.»

«Gut, hast du mit dem Mann einen Begleiter an deiner Seite. Das beruhigt
mich, Louis. Nach diesen schweren Wochen. Doch erstmal brauche ich Zeit,
das Ganze zu verdauen. Es ist ganz schön viel auf einmal.»

«Ja, ich weiss. Und Maman, Geheimnisse dieser Art sind vorbei. Jetzt
kümmere ich mich wirklich um die Dinge. Und sobald ich klarer bin, bist
du, wirst du die Erste sein, die ich bei meiner Rückkehr besuche.»

\sectionbreak

Eine willkommene Ruhe stellte sich ein, breitete sich gelassen aus und
während Louis seine Gitarre ansah, wie sie dastand, angelehnt an den
kleinen Tisch, bedingungslos wartend, erreichte ihn ein Funken Glück.

Er nahm das Instrument mit Ehrfurcht an sich. Und vielleicht das erste
Mal seit Wochen flogen seine Finger federleicht über die Saiten.

\qrblock{HAEVN \emph{«Where the Heart is»}}{media/QR-Codes/052_Where_the_Heart_is_Haevn_Spotify.png}
